% 对应题号 M2-035，本题条件、变量及几何关系逐项核对。
% 修订：启用中文断行；核对题设及图文对应。
\documentclass[tikz,border=14pt]{standalone}
\usepackage{fontspec,xeCJK}
\usepackage{amsmath,amssymb}
\usepackage{xcolor}
\usepackage{tikz}
\usetikzlibrary{arrows.meta,calc,patterns,positioning,bending}
\IfFontExistsTF{Heiti SC}{\IfFontExistsTF{Heiti SC}{\setmainfont{Heiti SC}}{\setmainfont{FandolHei-Regular}}\IfFontExistsTF{Heiti SC}{\setsansfont{Heiti SC}}{\setsansfont{TeX Gyre Heros}}}{\IfFontExistsTF{Songti SC}{\setmainfont{Songti SC}}{\setmainfont{FandolSong-Regular}}\IfFontExistsTF{Songti SC}{\setsansfont{Songti SC}}{\setsansfont{FandolSong-Regular}}}
\definecolor{bluefield}{RGB}{24,76,255}
\definecolor{redfield}{RGB}{232,49,35}
\definecolor{orangefield}{RGB}{236,132,18}
\definecolor{greenfield}{RGB}{18,132,86}
\definecolor{purplefield}{RGB}{112,72,190}
\definecolor{linegray}{RGB}{60,68,78}
\definecolor{softblue}{RGB}{238,243,255}
\definecolor{softgreen}{RGB}{238,249,244}
\definecolor{softorange}{RGB}{255,247,235}
\tikzset{
  arrow/.style={-{Latex[length=3.1mm,width=2.1mm]},line width=1.1pt},
  thinarr/.style={-{Latex[length=2.5mm,width=1.7mm]},line width=.85pt},
  axis/.style={-{Latex[length=2.5mm,width=1.7mm]},line width=.8pt,draw=linegray},
  guide/.style={draw=linegray!55,dashed,line width=.65pt},
  box/.style={draw=linegray,rounded corners=3pt,fill=white,align=center,inner sep=6pt,line width=.75pt},
  bluebox/.style={box,draw=bluefield,fill=softblue},
  greenbox/.style={box,draw=greenfield,fill=softgreen},
  orangebox/.style={box,draw=orangefield,fill=softorange}
}
\IfFontExistsTF{Heiti SC}{\setCJKmainfont[AutoFakeBold=1.5]{Heiti SC}}{\setCJKmainfont[AutoFakeBold=1.5]{FandolHei-Regular}}
\IfFontExistsTF{Heiti SC}{\setCJKsansfont{Heiti SC}}{\setCJKsansfont{FandolHei-Regular}}
\begin{document}

\begin{tikzpicture}[x=1cm,y=1cm]
\path[use as bounding box] (-10,-4.6) rectangle (10,4.6);
\fill[white] (-10,-4.6) rectangle (10,4.6);
\node[font=\Large\bfseries,text width=18.6cm,align=center] at (0,4.05) {平面四电荷：单极与偶极相消，四极张量记录方向差异};
\begin{scope}[shift={(-4.8,-.15)}]
\draw[axis] (-2.7,0)--(2.7,0) node[right] {$x$}; \draw[axis] (0,-2.7)--(0,2.7) node[above] {$y$};
\foreach \x in {-1.6,1.6}{\fill[redfield] (\x,0) circle(.12) node[above] {$+q$};}
\foreach \y in {-1.6,1.6}{\fill[bluefield] (0,\y) circle(.12) node[right] {$-q$};}
\draw[redfield,line width=1pt] (-2.1,-.5) rectangle (2.1,.5); \draw[bluefield,line width=1pt] (-.5,-2.1) rectangle (.5,2.1);
\end{scope}
\begin{scope}[shift={(4.65,-.1)}]
\node[bluebox,text width=6.8cm] at (0,1.65) {$Q=0,\qquad \mathbf p=\mathbf0$};
\node[greenbox,text width=6.8cm] at (0,0) {$\displaystyle \mathcal Q=6qa^2\,\mathrm{diag}(1,-1,0)$};
\node[orangebox,text width=6.8cm] at (0,-1.7) {$\displaystyle \varphi_2=\frac{k}{2r^3}\mathcal Q_{ij}\hat r_i\hat r_j=\frac{3kqa^2}{r^3}\sin^2\theta\cos2\phi$};
\end{scope}
\end{tikzpicture}
\end{document}
